\paragraph{Графический интерфейс.}
Графический интерфейс реализован в двух вариантах: настольное приложение на \verb|tkinter| в модуле \verb|app_gui.py| и веб-интерфейс (HTML/CSS/JavaScript + FastAPI). Веб-интерфейс включает:
\begin{itemize}
  \item вкладку «Документы корпуса»~--- просмотр списка документов с пагинацией, фильтрация по выбранным документам, удаление, просмотр полного текста по двойному щелчку, отображение суммы токенов по корпусу;
  \item вкладку «Частотные характеристики»~--- три режима (словоформы, леммы, части речи) с визуальным выделением выбранного режима, таблица частот с пагинацией, выделение строки леммы и кнопка «Скрыть лемму» для исключения леммы из учёта;
  \item вкладку «Конкорданс»~--- таблица контекстов вхождений запроса с пагинацией;
  \item вкладку «Морфология леммы»~--- таблица словоформ и частот для введённой леммы с пагинацией;
  \item вкладку «Скрытые леммы»~--- список исключённых из корпуса лемм с возможностью восстановления;
  \item вкладку «Метаданные»~--- вывод библиографической информации по выбранному документу;
  \item общую панель пагинации для вкладок с таблицами (документы, частоты, конкорданс, морфология);
  \item панель поиска с выбором типа запроса (по лемме / по словоформе) и размера контекста.
\end{itemize}
